% META: {"id":"TCOMPL-AIR-EX-003","chapitre":"TCOMPL-CALCULS-AIRES","type_objet":"exercice","capacites":["TCOMPL-AIR-C3"],"capacites_codes":["C3"],"methodes":[],"parcours":2,"competences":["calculer"],"duree_min":20,"mode_creation":"ex_nihilo","sources_inspiration":[],"fichier_tex":"chapitres/TCOMPL-CALCULS-AIRES/exercices/TCOMPL-AIR-EX-003.tex","corrige_tex":"chapitres/TCOMPL-CALCULS-AIRES/corriges/TCOMPL-AIR-CO-003.tex","status":"approved","origin":"generated","status_history":["generated","audited","accepted","approved"],"release_acceptance":"RELEASE_OWNER_BATCH_ACCEPTANCE","acceptance_closure_digest":"sha256:9b3ccf9a81c5520fbb7e03b7b2d3e7bf2057a02834b6d908bf3be5f49f3b553f","acceptance_receipt_digest":"sha256:4fad86f0282e0fa4e0419cca1ad6379ae7af5a280c04a4a90880f90a1c044242"}
% BEGIN-VERIFY
% from sympy import symbols, solve, Eq, integrate, Rational
% x = symbols('x')
% f = -x**2 + 4
% g = x**2 - 2*x
% inter = solve(Eq(f, g), x)
% assert inter == [-1, 2]
% aire = integrate(f - g, (x, -1, 2))
% assert aire == Rational(9, 1)
% END-VERIFY
\begin{exercice}{TCOMPL-AIR-EX-003}{2}{20}
Soit $f(x) = -x^2+4$ et $g(x) = x^2-2x$.
\begin{enumerate}
  \item Déterminer les abscisses des points d'intersection des courbes de $f$ et $g$.
  \item Justifier que $f \geq g$ entre ces deux abscisses.
  \item Calculer l'aire comprise entre les deux courbes.
\end{enumerate}
\end{exercice}
